\documentclass[11pt,a4paper]{article}
\usepackage[T1]{fontenc}
\usepackage[utf8]{inputenc}
\usepackage{graphicx}
\usepackage{booktabs}
\usepackage{amsmath}
\usepackage{amssymb}
\usepackage{hyperref}
\usepackage[margin=1in]{geometry}
\usepackage{float}
\usepackage{longtable}
\usepackage{multirow}
\usepackage{array}

\graphicspath{{../figures/}{figures/}}

\title{\textbf{GlycoSMILES2BAP: An Automated Pipeline for Stereochemistry-Preserving Glycan Structure Prediction with AlphaFold 3}}

\author{Qiang Xia\\
\small Zhejiang Xinghe Tea Technology Co., Ltd.\\
\small Hangzhou, Zhejiang, China\\
\small \texttt{xiaqiang@xinghetea.com}}

\date{}

\begin{document}

\maketitle

\begin{abstract}
\noindent\textbf{Motivation:} AlphaFold 3 (AF3) has revolutionized protein structure prediction and expanded capabilities to include glycan ligands. However, recent work by Huang et al.\ (2025) identified a critical limitation: standard input formats (SMILES, userCCD) produce stereochemically incorrect glycan structures, while the stereochemistry-preserving CCD+bondedAtomPairs (BAP) format requires impractical manual atom-by-atom specification.

\noindent\textbf{Results:} We present GlycoSMILES2BAP, an automated pipeline that converts standard glycan notations (IUPAC-condensed, WURCS) to AF3-compatible CCD+BAP format. The pipeline integrates three core modules: (1) a CCD mapper supporting 28+ monosaccharide configurations, (2) a topology parser extracting linkage information, and (3) a BAP generator producing explicit atom-pair bond specifications. Validated on a curated benchmark of 50 diverse glycan structures with systematic evaluation, GlycoSMILES2BAP achieves 97.8\% epimer accuracy, 97.4\% anomeric accuracy, and 95.9\% linkage accuracy compared to ${\sim}60$\% for SMILES-based approaches, with processing time $<$1 second per structure versus 30--60 minutes for manual specification. Systematic ablation studies confirm that all three core modules contribute significantly to the high accuracy. Furthermore, validation against 10 literature-reported structure errors demonstrated 100\% correction rate, and database-scale processing demonstrated scalability with 94\% success rate on 100 GlyTouCan representative structures (average processing time 0.82ms per structure).

\noindent\textbf{Conclusions:} GlycoSMILES2BAP eliminates the input format barrier for AF3 glycan modeling, enabling accurate, reproducible structure prediction for the glycobiology community. The tool successfully corrects common stereochemistry errors and scales to database-level processing, making it suitable for both individual structure validation and large-scale structural glycomics. The tool is open-source and readily integrable with existing glycan databases (GlyTouCan, GlyGen).

\noindent\textbf{Availability:} Open-source implementation available at \url{https://github.com/xiaqiang/glycosmiles2bap}

\noindent\textbf{Keywords:} AlphaFold 3, glycans, stereochemistry, structure prediction, CCD, bondedAtomPairs
\end{abstract}

\section{Introduction}

Glycans are essential biological molecules involved in protein folding, cell signaling, immune recognition, and pathogen interaction. Over 50\% of human proteins undergo glycosylation, making accurate glycan structure prediction crucial for understanding biological mechanisms and developing therapeutics. GlyTouCan, the international glycan structure repository, catalogs over 200,000 unique glycan structures, highlighting the scale of structural diversity that researchers need to navigate. This diversity---arising from variations in monosaccharide composition, linkage positions, anomeric configurations, and branching patterns---poses unique challenges for computational structure prediction.

AlphaFold 3 (AF3) represents a breakthrough in biomolecular structure prediction, achieving unprecedented accuracy for protein-ligand complexes including glycans. This advancement has generated considerable excitement in the glycobiology community, as accurate glycan structure prediction could accelerate research in glycoprotein engineering, vaccine design, and therapeutic glycosylation.

\subsection{The Stereochemistry Problem in AF3 Glycan Modeling}

However, a fundamental challenge emerges when using AF3